\noindent
{\Large \textbf{\section{Постановка задачи}}}\par

Реализовать двумерный клеточный автомат с окрестностью фон Неймана в соответствии с полученным номером $N$. Граничные условия - тороидальные. Пользователь определяет ширину поля и количество итераций. Необходимо учесть возможность ввода различных начальных условий (как вручную, так и случайным образом) по выбору пользователя.

Номер вычисляется по формуле
\[
 N = \text{номер варианта} \cdot 11 \cdot \text{год рождения} \cdot \text{день} \cdot \text{месяц}.
\]
Для номера варианта $=9$: $N = 9 \cdot 11 \cdot 2006 \cdot 21 \cdot 12 = 50045688$.

\newpage
